\documentclass{amsart}
% For dealing with references we use the comment environment
\usepackage{verbatim}
\newenvironment{reference}{\comment}{\endcomment}
%\newenvironment{reference}{}{}
\newenvironment{slogan}{\comment}{\endcomment}
\newenvironment{history}{\comment}{\endcomment}

% For commutative diagrams we use Xy-pic
\usepackage[all]{xy}

% We use 2cell for 2-commutative diagrams.
\xyoption{2cell}
\UseAllTwocells

% We use multicol for the list of chapters between chapters
\usepackage{multicol}

% This is generally recommended for better output
\usepackage{lmodern}
\usepackage[T1]{fontenc}

% For cross-file-references

% Package for hypertext links:
\usepackage{hyperref}

% For any local file, say "hello.tex" you want to link to please

% Theorem environments.
%
\theoremstyle{plain}
\newtheorem{theorem}[subsection]{Theorem}
\newtheorem{proposition}[subsection]{Proposition}
\newtheorem{lemma}[subsection]{Lemma}

\theoremstyle{definition}
\newtheorem{definition}[subsection]{Definition}
\newtheorem{example}[subsection]{Example}
\newtheorem{exercise}[subsection]{Exercise}
\newtheorem{situation}[subsection]{Situation}

\theoremstyle{remark}
\newtheorem{remark}[subsection]{Remark}
\newtheorem{remarks}[subsection]{Remarks}

\numberwithin{equation}{subsection}

% Macros
%
\def\lim{\mathop{\mathrm{lim}}\nolimits}
\def\colim{\mathop{\mathrm{colim}}\nolimits}
\def\Spec{\mathop{\mathrm{Spec}}}
\def\Hom{\mathop{\mathrm{Hom}}\nolimits}
\def\Ext{\mathop{\mathrm{Ext}}\nolimits}
\def\SheafHom{\mathop{\mathcal{H}\!\mathit{om}}\nolimits}
\def\SheafExt{\mathop{\mathcal{E}\!\mathit{xt}}\nolimits}
\def\Sch{\mathit{Sch}}
\def\Mor{\mathop{\mathrm{Mor}}\nolimits}
\def\Ob{\mathop{\mathrm{Ob}}\nolimits}
\def\Sh{\mathop{\mathit{Sh}}\nolimits}
\def\NL{\mathop{N\!L}\nolimits}
\def\CH{\mathop{\mathrm{CH}}\nolimits}
\def\proetale{{pro\text{-}\acute{e}tale}}
\def\etale{{\acute{e}tale}}
\def\QCoh{\mathit{QCoh}}
\def\Ker{\mathop{\mathrm{Ker}}}
\def\Im{\mathop{\mathrm{Im}}}
\def\Coker{\mathop{\mathrm{Coker}}}
\def\Coim{\mathop{\mathrm{Coim}}}

% Boxtimes
%
\DeclareMathSymbol{\boxtimes}{\mathbin}{AMSa}{"02}

%
% Macros for moduli stacks/spaces
%
\def\QCohstack{\mathcal{QC}\!\mathit{oh}}
\def\Cohstack{\mathcal{C}\!\mathit{oh}}
\def\Spacesstack{\mathcal{S}\!\mathit{paces}}
\def\Quotfunctor{\mathrm{Quot}}
\def\Hilbfunctor{\mathrm{Hilb}}
\def\Curvesstack{\mathcal{C}\!\mathit{urves}}
\def\Polarizedstack{\mathcal{P}\!\mathit{olarized}}
\def\Complexesstack{\mathcal{C}\!\mathit{omplexes}}
% \Pic is the operator that assigns to X its picard group, usage \Pic(X)
% \Picardstack_{X/B} denotes the Picard stack of X over B
% \Picardfunctor_{X/B} denotes the Picard functor of X over B
\def\Pic{\mathop{\mathrm{Pic}}\nolimits}
\def\Picardstack{\mathcal{P}\!\mathit{ic}}
\def\Picardfunctor{\mathrm{Pic}}
\def\Deformationcategory{\mathcal{D}\!\mathit{ef}}

\setlength{\emergencystretch}{3em}
\begin{document}
\title[Stacks Project, Tag 0H0U]{291 --- Unbounded quasi-coherent complexes satisfy affine fppf cohomological descent}
\maketitle
\noindent Copyright (C) 2005--2025 Johan de Jong. Stacks Project authors. Source: \href{https://stacks.math.columbia.edu/tag/0H0U}{Tag 0H0U}.

\noindent Editorial difficulty note (not author text). Difficulty H2: The bounded-below comparison does not automatically apply to arbitrary unbounded complexes. The proof must reconstruct the sheaf complex as a derived inverse limit of truncations using eventual stabilization on every affine object, and justify passing derived global sections through that inverse limit..

\noindent Editorial provenance note (not author text). History: excerpt from revision \texttt{a0444\allowbreak 6e57e\allowbreak c1fbc\allowbreak 252a8\allowbreak 71afc\allowbreak ec775\allowbreak 2fb28\allowbreak 07b14}, etale-cohomology.tex, lines 22046--22122. Reference presentation was adapted; original-author context and core support blocks were retained or added. The mathematical wording is unchanged. Licensed under GNU FDL 1.2 or later; no invariant sections or cover texts. See ../licenses/COPYING. Context blocks reproduce author definitions and supporting results; they are not additional counted problems.

\begin{lemma}
\label{lemma-cohomological-descent-complex-modules}
Let $S = \Spec(A)$ be an affine scheme. Let $M^\bullet$ be a complex
of $A$-modules. Consider the complex $\mathcal{F}^\bullet$ of
presheaves of $\mathcal{O}$-modules on
$(\textit{Aff}/S)_{fppf}$ given by the rule
$$
(U/S) = (\Spec(B)/\Spec(A)) \longmapsto M^\bullet \otimes_A B
$$
Then this is a complex of modules and the canonical map
$$
M^\bullet \longrightarrow
R\Gamma((\textit{Aff}/S)_{fppf}, \mathcal{F}^\bullet)
$$
is a quasi-isomorphism.
\end{lemma}

\begin{proof}
Each $\mathcal{F}^n$ is a sheaf of modules as it agrees with the
restriction of the module $\mathcal{G}^n = (\widetilde{M}^n)^a$ of
Lemma \href{https://stacks.math.columbia.edu/tag/0DEW}{lemma review quasi coherent (Tag 0DEW)} to
$(\textit{Aff}/S)_{fppf} \subset (\Sch/S)_{fppf}$.
Since this inclusion defines an equivalence of ringed topoi
(Topologies, Lemma \href{https://stacks.math.columbia.edu/tag/021V}{lemma affine big site fppf (Tag 021V)}),
we have
$$
R\Gamma((\textit{Aff}/S)_{fppf}, \mathcal{F}^\bullet) =
R\Gamma((\Sch/S)_{fppf}, \mathcal{G}^\bullet)
$$
We observe that $M^\bullet = R\Gamma(S, \widetilde{M}^\bullet)$
for example by Derived Categories of Schemes, Lemma
\href{https://stacks.math.columbia.edu/tag/06Z0}{lemma affine compare bounded (Tag 06Z0)}. Hence we are trying
to show the comparison map
$$
R\Gamma(S, \widetilde{M}^\bullet)
\longrightarrow
R\Gamma((\Sch/S)_{fppf}, (\widetilde{M}^\bullet)^a)
$$
is an isomorphism. If $M^\bullet$ is bounded below, then this
holds by Descent, Proposition
\href{https://stacks.math.columbia.edu/tag/03DW}{proposition same cohomology quasi coherent (Tag 03DW)}
and the first spectral sequence of Derived Categories, Lemma
\href{https://stacks.math.columbia.edu/tag/015J}{lemma two ss complex functor (Tag 015J)}.
For the general case, let us write $M^\bullet = \lim M_n^\bullet$
with $M_n^\bullet = \tau_{\geq -n}M^\bullet$. Whence the system
$M_n^p$ is eventually constant with value $M^p$. We claim that
$$
(\widetilde{M}^\bullet)^a = R\lim (\widetilde{M}_n^\bullet)^a
$$
Namely, it suffices to show that the natural map from left to right
induces an isomorphism on cohomology over any affine object
$U = \Spec(B)$ of $(\Sch/S)_{fppf}$. For $i \in \mathbf{Z}$ and
$n > |i|$ we have
$$
H^i(U, (\widetilde{M}_n^\bullet)^a) =
H^i(\tau_{\geq -n}M^\bullet \otimes_A B) =
H^i(M^\bullet \otimes_A B)
$$
The first equality holds by the bounded below case treated above.
Thus we see from Cohomology on Sites, Lemma
\href{https://stacks.math.columbia.edu/tag/0D6K}{lemma RGamma commutes with Rlim (Tag 0D6K)}
that the claim holds. Then we finally get
\begin{align*}
R\Gamma((\Sch/S)_{fppf}, (\widetilde{M}^\bullet)^a)
& =
R\Gamma((\Sch/S)_{fppf}, R\lim (\widetilde{M}_n^\bullet)^a) \\
& =
R\lim R\Gamma((\Sch/S)_{fppf}, (\widetilde{M}_n^\bullet)^a) \\
& =
R\lim M_n^\bullet \\
& =
M^\bullet
\end{align*}
as desired. The second equality holds because $R\lim$ commutes with
$R\Gamma$, see Cohomology on Sites, Lemma
\href{https://stacks.math.columbia.edu/tag/0D6K}{lemma RGamma commutes with Rlim (Tag 0D6K)}.
\end{proof}
\end{document}
